%start% Brought this in from Patrick thesis
\documentclass{thesis}
\usepackage[utf8]{inputenc}
\usepackage{parskip}
\usepackage{hepnames}
\usepackage[printonlyused]{acronym}
\usepackage{thesis}
\usepackage{ptdr-definitions}
\usepackage{hepnicenames}
\usepackage{hepunits}
\usepackage{abhep}
\usepackage{slashed}
\usepackage{multirow}
%end%

%\documentclass[hyperpdf]{hepthesis}
%\documentclass[a4paper,10pt]{book}
%\usepackage{fullpage}
%\usepackage[margin=2cm]{geometry} % Defining geometry of the usable text space

%% Using Babel allows other languages to be used and mixed-in easily
%\usepackage[ngerman,english]{babel}
% \usepackage[english]{babel}
% \selectlanguage{english}

%% Citation system tweaks
% \usepackage{cite}

% \usepackage{amsmath}
% \usepackage{amsthm}
% \usepackage{amsfonts}
% \usepackage{amssymb}
% \usepackage{graphicx}
% \usepackage{color}                 % For color (highlights,...) 
% \usepackage{caption}               % ?????
% \usepackage{soul}                  % For text highlights
% \usepackage{threeparttable}        % Better tables
% \usepackage[xindy,toc,acronym]{glossaries} %(recommended) for the indexing phase, as opposed to makeindex (the default)
% \usepackage{hyperref}              % Need after glossaries (allows links)

% % \newacronym[longplural={Frames per Second}]{fpsLabel}{FPS}{Frame per Second}

%A
\newacronym{ALICE}{ALICE}{A Large Ion Collider Experiment}
\newacronym{ATLAS}{ATLAS}{A Toroidal LHC ApparatuS}
\newacronym{APD}  {APD}  {Avalanche photo-diodes}

%B
\newacronym{BSM}{BSM}{Beyond the Standard Model}
\newacronym{BDT}{BDT}{Boosted Decision Tree}

%C
\newacronym{CERN}{CERN}{European Organization for Nuclear Research} %derived from the name Conseil Européen pour la Recherche Nucléaire
\newacronym{CMS} {CMS} {Compact Muon Solenoid}

%D
\newacronym{DAQ}{DAQ}{Data Acquisition}
\newacronym{DQM}{DQM}{Data Quality Monitoring}

%E
\newacronym{EB}  {EB}  {ECAL Barrel}
\newacronym{EE}  {EE}  {ECAL Endcap}
\newacronym{ECAL}{ECAL}{Electromagnetic Calorimeter}

%F
\newacronym{FCT}{FCT}   {Fundação para a Ciência e a Tecnologia}

%G

%H
\newacronym{HCAL}{HCAL}  {Hadronic Calorimeter}
\newacronym{HLT} {HLT}   {High Level Trigger}

%I

%J
\newacronym{JER}{JER}{jet energy resolution}
\newacronym{JES}{JES}{jet energy scale}

%K

%L
\newacronym{L1T}   {L1T}   {Level 1 Trigger}
\newacronym{LINAC2}{LINAC2}{Linear Particle Accelerator 2}
\newacronym{LHC}   {LHC}   {Large Hadron Collider}

%M
\newacronym{MC} {MC} {Monte Carlo}
\newacronym{MET}{MET}{Missing Transverse Energy}

%N

%O

%P
\newacronym{PS} {PS} {Proton Synchrotron}
\newacronym{PSB}{PSB}{Proton Synchrotron Booster}

%Q

%R

%S
\newacronym{SM} {SM} {Standard Model}
\newacronym{SPS}{SPS}{Super  Proton Synchrotron}

%T

%U

%V
\newacronym{VBF}{VBF}{vector Boson Fusion}
\newacronym{VPT}{VPT}{Vacuum Photo-Triodes}
%W

%X

%Y

%Z


% \acro{CL}{confidence level}
% \acro{CSC}{cathode strip chamber}
% \acro{CSV}{combined secondary vertex}
% \acro{CTF}{combinatorial track finder}
% \acro{DA}{deterministic annealing}
% \acro{DT}{drift tube}
% \acro{GSF}{Gaussian sum filter}
% \acro{HB}{hadron barrel}
% \acro{HE}{hadron endcaps}
% \acro{HF}{hadron forward}
% \acro{HLT}{high-level trigger}
% \acro{HO}{hadron outer}
% \acro{HPS}{hadron plus strips}
% \acro{JER}{jet energy resolution}
% \acro{JES}{jet energy scale}
% \acro{L1}{Level-1}
% \acro{LHCHXSWG}{LHC Higgs Cross Section Working Group}
% \acro{LO}{leading order}
% \acro{MPF}{missing transverse energy projection fraction}
% \acro{MPI}{multi-parton interaction}
% \acro{MVA}{multi-variate analysis}
% \acro{MSSM}{minimal supersymmetric standard model}
% \acro{NLO}{next-to-leading order}
% \acro{NNLO}{next-to-next-to-leading order}
% \acro{NNLL}{next-to-next-to-leading logarithmic}
% \acro{pdf}{probability density function}
% \acro{PDF}{parton distribution function}
% \acro{PF}{particle flow}
% \acro{PS}{Proton Synchrotron}
% \acro{PSB}{Proton Synchrotron Booster}
% \acro{PU}{pile-up}
% \acro{PV}{Primary Vertex}
% \acro{QCD}{Quantum Chromodynamics}
% \acro{RF}{radio frequency}
% \acro{RPC}{resistive plate chamber}
% \acro{SPS}{Super Proton Synchrotron}
% \acro{SSV}{simple secondary vertex}
% \acro{SM}{standard model}
% \acro{TEC}{tracker endcaps}
% \acro{TIB}{tracker inner barrel}
% \acro{TID}{tracker inner disks}
% \acro{TOB}{tracker outer barrel}
% \acro{UE}{underlying event}
% \acro{VBF}{vector boson fusion}
% \acro{WLCG}{Worldwide LHC Computing Grid}


\title{Search for Higgs Decay to Dark Matter and Trigger Studies} 
\author{João Carlos Arnauth Pela}
% \email[Contact email: ]{joaopela@gmail.com}
% \affiliation{Imperial College London}
% \keywords{Particle Physics}
\date{21 September 2015}

%% Doc-specific PDF metadata
\makeatletter
\@ifpackageloaded{hyperref}{
\hypersetup{
  pdftitle    = {Search for Higgs Decay to Dark Matter and Trigger Studies},
  pdfsubject  = {João Pela's PhD thesis},
  pdfkeywords = {LHC, CMS, Higgs, Dark Matter, Trigger},
  pdfauthor   = {\textcopyright\ João Pela}
}
}{}
\makeatother

\begin{document}

\begin{frontmatter}
  %% Title
%!!UPDATE THIS
\titlepage[of Imperial College London]%
{A dissertation submitted to Imperial College London\\
  for the degree of Doctor of Philosophy}
  
\clearpage
\textbf{The copyright of this thesis rests with the author and is made available under a Creative Commons 
Attribution Non-Commercial No Derivatives licence. Researchers are free to copy, distribute or 
transmit the thesis on the condition that they attribute it, that they do not use it for commercial 
purposes and that they do not alter, transform or build upon it. For any reuse or redistribution, 
researchers must make clear to others the licence terms of this work.}


%% Abstract
\begin{abstract}%[\smaller \thetitle\\ \vspace*{1cm} \smaller {\theauthor}]
The Compact Muon Solenoid (CMS) is a general-purpose particle detector at the CERN Large Hadron Collider (LHC). The goal of this experiment is to search for the Higgs boson and evidence of new physics and to test the prediction of the Standard Model (SM) at the TeV scale. This thesis describes the analysis of proton-proton collision data recorded by CMS during 2012 and support work for data taking during the same period.

A search for the Vector Boson Fusion (VBF) produced Higgs boson invisible decays, using $19.5\,\femto\barn^-1$ of data recorded with prompt reconstruction triggers at a center of mass energy of $8\,\TeV$, is presented. Events are selected with two forward jets and large Missing Transverse Momentum. Assuming the SM VBF production cross section and acceptance, the observed (expected) upper limit at the 95\% confidence level on the \BRinv\, is determined to be of 65\% (49\%) for $m_H=125\,\GeV$.

A second search for the VBF Higgs boson invisible decays, using $19.2\,\femto\barn^-1$ of data recorded with delayed reconstruction (parked) triggers at a center of mass energy of $8\,\TeV$, is also presented. A new event selection criteria was developed taking advantage of the lower trigger requirements. Assuming the SM VBF Higgs production cross section and acceptance, the observed (expected) upper limit at the 95\% confidence level on the \BRinv\, is determined to be of 57\% (40\%) for $m_H=125\,\GeV$.

Monitoring for the CMS Level 1 Trigger system has been developed and used during the 2012 and subsequent LHC data acquisition periods. Contribution to the high reliability of this system during data taking and providing crucial information for validation of the data quality.
\end{abstract}

%DONE

%% Declaration
\begin{declaration}
I declare that the work contained in this thesis is my own, and all results and figures taken from other sources are indicated in the text and referenced appropriately. The analyses presented in this thesis were developed in close collaboration with other members of the CMS experiment.

For the Vector Boson Fusion (VBF) Higgs to invisible prompt analysis I have contributed with the literature review and final background cross section input values for background normalization. QCD multi-jet background studies were also preformed with the target of improving the final selection or prepare the parked analysis, this analysis is presented to give context for the more important work developed later. 

For the VBF Higgs to invisible parked analysis I have participated in the development of the parked trigger, which was used to record the majority of the analysed data. I have continued the QCD multi-jet background studies which have lead to the production of dedicated Monte Carlo simulations and novel approaches to reject this type of events. I have been the responsible for the cross check analysis of the main result which has successfully validated the implementation of the main analysis \cite{ARTICLE:CMSVBFHiggsInvisibleParkedAnalysisPAS}. It is a normal a requirement for many \gls{CMS} publications to have a cross check analysis implemented independently from the main result in order to be able to ensure the accuracy of the final results. I have also participated in the preparation of the Run II analysis where I have lead the development of both the triggers used for data recording during 2015. Additionally, I have developed a method to create the first QCD multi-jet Monte Carlo sample with no MET requirements with signal like properties.

As part of the CMS Level 1 Trigger (L1T) Detector Performance Group (DPG) I have developed monitoring tools for this system, which were both used for real-time monitoring and posterior data certification for physics analysis usage. My work in this group has lead to my appointment for two years to the position of coordinator of the \gls{CMS} \gls{L1T} Data quality Monitoring software development team. Field work was also performed by doing shifts as Trigger and Shift Leader in the experiment control room and on call shifts as the Trigger Detector on Call (DOC) expert.
 
\vspace*{1cm}
\begin{flushright}
João Pela
\end{flushright}
\end{declaration}

%% Acknowledgements
\begin{acknowledgements}
I would like to thank the Imperial College HEP Group and the FCT for supporting and giving me the opportunity to pursue research including my time spent at CERN. Thanks to my supervisor, David Colling, for his guidance during the years of my PhD, the many useful discussions and his enthusiasm for particle physics research. Thanks must also go to all my colleagues of the CMS VBF Higgs to invisible group with whom I have worked directly, Patrick Dunne, Anne-Marie Magnan, Alexandre Nikitenko, Jim Brooke and Chayanit Asawatangtrakuldee. Thanks to everyone in CMS L1T DPG group, specially to Carlo Battilana and Arno Heister for their support. Thanks to all my friends that supported me on my way to finish this degree, it is difficult to mention all of you, but here is a small list, Antonio Pacheco, André David, Alice Mitchell, Chris Holdsworth-Swan, Francisco Sousa, Glenn Spiers, Guillermo Marrero, João Filipe, Miguel Cunhal, Miguel Machado, Nadia Silva, Raquel Monteiro, and many, many more. Finally, and most importantly, I thank my parents and grandparents for their constant support, love and encouragement.
\end{acknowledgements}

\clearpage

\vspace*{\fill}
The work presented in this thesis was supported by the Portuguese Government through \gls{FCT} in the form of my PhD grant with the reference SFRH/BD/77979/2011. I am thankful for their support which allowed me to attain higher education.
\vspace*{\fill}

\begin{center}
\resizebox{1.0\linewidth}{!}{
\includegraphics{FrontMatter/Images/FundingBand.png}
}
\end{center}


% \begin{preface}
% Thesis structure and so on...
% \end{preface}

\dedication{To my grandmother.}

%% Preface
%\begin{preface}
%\end{preface}

%% ToC
\tableofcontents

\newpage
\listoffigures

\newpage
\listoftables

%% Strictly optional!
%\frontquote%
%{Writing in English is the most ingenious torture\\
%   ever devised for sins committed in previous lives.}%
%  {James Joyce}
 
% \titlepage[of the High Energy Physics Group]{%
%   A dissertation submitted to the Imperial College London\\ for the degree of Doctor of Philosophy}

\end{frontmatter}

\begin{mainmatter}
  \chapter{Theory}
%NOTE: Maybe would be good to make a more descriptive title

\colorbox{red}{TODO: Something}

\section{Standard Model of Particle Physics}

\colorbox{red}{TODO: Very brief summary of the Standard Model.}

% The \gls{sm} of particle physics is the currently accepted model for describing the physics of elementary particles.

\section{Higgs Mechanism}

Summary of the Higgs Mechanism. Should include
\begin{itemize}
 \item Motivations 
 \item Explanation of the mechanism itself
 \item Consequences 
 \item Possible decays
\end{itemize}

\section{Higgs Invisible decays}

\colorbox{red}{TODO: Explain what are SM Higgs invisible decays. Go over the possibility of BSM invisible decays.}

  \chapter{Experimental Apparatus}
\label{CHAPTER:ExperimentalApparatus}




% To cite
The LHC Machine article \cite{ARTICLE:LHC Machine}.
%

\section{The Large Hadron Collider}
\label{SECTION:ExperimentalApparatus_LHC}

\colorbox{red}{
\begin{minipage}{\linewidth}
TODO: 

\begin{itemize}
  \item CERN basics location etc
  \item Basics of machine and operation
  \item New to include Instantaneous luminosity equation
\end{itemize}

\end{minipage}
}

The Large Hadron Collider \cite{ARTICLE:LHCMachine} is a 27 km synchrotron machine located in Geneva Switzerland.

Luminosity Equation

\begin{equation}
L=\frac{N_{b}^{2}n_{b}f_{\text{rev}}\gamma}{4\pi\epsilon_{n}\beta^{*}}F,
\end{equation}

\section{The Compact Muon Solenoid Experiment}
\label{SECTION:ExperimentalApparatus_CMS}





\subsection{Tracker}
\label{SUBSECTION:ExperimentalApparatus_CMS_Tracker}

\subsection{Electromagnetic Calorimeter}
\label{SUBSECTION:ExperimentalApparatus_CMS_ECAL}

\subsection{Hadronic Calorimeter}
\label{SUBSECTION:ExperimentalApparatus_CMS_HCAL}

\subsection{Solenoid Magnet}
\label{SUBSECTION:ExperimentalApparatus_CMS_Magnet}

\subsection{Muon System}
\label{SUBSECTION:ExperimentalApparatus_CMS_Mouns}

\subsection{Data Acquisition System}
\label{SUBSECTION:ExperimentalApparatus_CMS_DAQ}

\subsection{Trigger System}
\label{SUBSECTION:ExperimentalApparatus_CMS_Trigger}

The upgrade tdr \cite{CMSL1UpgradeTDR}.

\subsection{Computing}
\label{SUBSECTION:ExperimentalApparatus_CMS_Computing}

% \gls{dqm} and \gls{dqm} 

\subsection{Run II Updgrades}
\label{SUBSECTION:ExperimentalApparatus_CMS_RUNII}


  \chapter{Physics Objects and Monte Carlo simulation}

\section{Physics objects definition}

\subsection{Electron}

\subsection{Muon}

\subsection{Tau}

\subsection{Jets}

\subsection{Missing Transverse Energy}

\section{Monte Carlo simulation}



% \clearpage
% \section{Introduction}

The search for an invisible decay of a vector boson produced Higgs boson was first made public with CMS Physics Analysis Summary (PAS) HIG-13-013 which was further improved and combined with other Higgs boson production channel in the CMS paper HIG-13-30. Additional support material can be found at the CMS Analysis Notes (AN) AN-2012/403\cite{CMS_AN_2013-403} and AN-2013/205.

During the 2012 data taking run two main streams of data were recorded. The main stream with an event rate of the order of 300 Hz to be promptly reconstructed and made available for analysis in a few days after being recorded, this dataset is referred to as the prompt data. The secondary stream with lower trigger thresholds with an event rate of the order of 1kHz which would only be reconstructed when the computing resources would be available outside of the data taking period, dataset is referred to as parked data. Our previous results
were produced using the prompt data only and this work now extends on previous work by using the now available parked data. Since this dataset has been recorded with lower trigger thresholds the analysis was re-optimised to take advantage of this new available phase space. The details of the newly developed analysis can be found in CMS AN-14-243\cite{CMS_AN_2014-243}.

It is normally a requirement for many CMS publications to have a cross check analysis implemented independently from the main result in order to be able to ensure accuracy of the final results due to possible errors with the software implementation. For this purpose the previous prompt data VBF Higgs to Invisible results and publication were produced by two different and independent code frameworks and before publication a good level of synchronization were obtained. Due to lack of man power and time it was decide for the 2012 parked data analysis to only proceed with a single framework. At a later stage of the analysis it was thought that at least some level of cross check would be a good measure to limit the possibility of implementation errors and to allow extra confidence on the final results.

This cross check analysis starts from the same ntuples produced by the main analysis which were produced over all the relevant datasets and are recorded with data formats also used by other analysis at Imperial College London, e.g. both the SM and MSSM Higgs to $\tau\bar{\tau}$, the Higgs to $\tau\bar{\tau}b\bar{b}$ and prompt Higgs to invisible analyses. No cuts are applied at ntuple production except the official CMS selection for good usable data using the appropriate golden JSON file.

From those initial ntuple an independent code framework was developed in order to replicate all relevant numbers and plots from the main analysis.

% \section{Data samples}

For this analysis we used the proton-proton collision parked data collected by the CMS experiment during 2012-13. An analysis purposely constructed trigger was used to collect this data.

\begin{table}[htp]
\centering

\begin{tabular}{|l|c|}
\hline
Dataset & $\int{Luminosity}$ $[pb^{-1}]$ \\
\hline \hline
/MET/Run2012A-22Jan2013-v1/AOD        & 889 \\
/VBF1Parked/Run2012B-22Jan2013-v1/AOD & 3871 \\
/VBF1Parked/Run2012C-22Jan2013-v1/AOD & 7152 \\
/VBF1Parked/Run2012D-22Jan2013-v1/AOD & 7317 \\
\hline
Total analysed & 19229 \\
\hline \hline
Total certified luminosity & 19789 \\
\hline
\end{tabular}

\caption{Used datasets and their respective integrated luminosities after applying certified for physics filtering. Each dataset corresponds to a recording period/era of the 2012-13 data acquisition run.}
\end{table}

% Review text abouve

% To include
% * Explain breifly the validadion process (JSON)
% * Table with data included


% \section{Event Filters and object definitions}

%STATUS: Finished
%TODO: 
% * Could not find any reference to this recommendations!
\subsection{Vertex}

The interaction point is normally assumed to be the reconstructed  primary vertex, defined as the vertex with highest sum of associated tracks \pt squared, or if that cannot be determined the beam spot position is assumed. Knowing precisely the interaction point will allow to determine object quantities relative to it which allow for better object identification and pile-up control. At least one good vertex implicitly required by tracking failure event quality filter. And Additionally we require explicitly that a good vertex is reconstructed with the following characteristics. 

\begin{itemize}
  \item NOT(isFake): We require a real reconstructed vertex from tracks, not the beam spot.
  \item Number of degrees of freedom: $n_{dof}>4$
  \item Longitudinal distance: $|z|<=24$
  \item Radial distance to beam line: $d_{xy}<2$
\end{itemize}

%STATUS: Finished
%TODO: 
% * Add citation for the Jet-MET POG for this filter usage (include data and page version)
% * QUESTION: Should I include description of what each filter does?
\subsection{Event quality filters}

During data recording issues may happen with the detector or acquisition which may render some of the events unusable. The groups responsible for each part of the detector and physics object check the data after and was taken and if they find such problems software event filters are made available for analysis to be able to remove this problematic events. This event filters will address issues like detector know problems, miss firing of calibration sequences or failure to reconstruct physics objects. The Jet-MET Particle Object Group (POG) recommends the usage of the following filters which are used in this analysis\cite{CMS:JetMETPOG:MissingETOptionalFilters}.

\begin{itemize}
  \item CSCTightHaloFilter
  \item HBHENoiseFilter
  \item EcalDeadCellTriggerPrimitiveFilter
  \item trackingFailureFilter
  \item eeBadScFilter
  \item ECAL Laser filter (via event list)
  \item HCAL Laser filter (via event list)
\end{itemize}

In turn the JetMET group recommend the usage of the following Tracking POG Filter\cite{CMS:TrackingPOG:TrackingPOGFilters}:

\begin{itemize}
  \item  logErrorTooManyClusters
  \item  manystripclus53X
  \item  toomanystripclus53X
\end{itemize}

The event rejection efficiencies from the all filters except ECAL and HCAL laser event filters can be found at table \ref{table_EventQualityFilterEff}. This values are measured the vertex requirements. 

\begin{table}[!htp]
\centering

\begin{tabular}{|l||c|c|c|c|}
\hline
Filter                  & Prompt A & Parked B & Parked C & Parked D \\
\hline \hline
ECAL Laser Filter       & 0.928521 & 0.008659 & 0.000000 & 0.000000 \\
HCAL Laser Filter       & 0.007258 & 0.000000 & 0.000270 & 0.000000 \\
\hline
ECAL+ HCAL Laser Filter & 0.935704 & 0.008659 & 0.000270 & 0.000000 \\
\hline
\end{tabular}

\caption{Event rejection efficiency for the ECAL and HCAL Laser Filters. This events have to be removed due to the untimely firing of the calibration laser for this systems.}
\label{table_CaloLaserFilterEff}
\end{table}


The values for the ECAL and HCAL laser event filters are present in table \ref{table_CaloLaserFilterEff}. This value are calculated after vertex requirement and the filters in table \ref{table_EventQualityFilterEff}. 

\begin{table}[htp]
\centering

\begin{tabular}{|l||c|c|c|c|}
\hline
Filter                             & Prompt A  & Parked B & Parked C & Parked D \\
\hline \hline
HBHENoiseFilter                    & 22.900905 & 0.190670 & 0.187739 & 0.170753 \\
EcalDeadCellTriggerPrimitiveFilter &  0.375381 & 0.009300 & 0.010206 & 0.012526 \\
eeBadScFilter                      &  0.007852 & 0.000001 & 0.000000 & 0.000009 \\
trackingFailureFilter              &  3.073876 & 0.000328 & 0.007464 & 0.000290 \\
manystripclus53X                   &  0.001829 & 0.001319 & 0.002335 & 0.001327 \\
toomanystripclus53X                &  0.000484 & 0.001149 & 0.002006 & 0.001173 \\
logErrorTooManyClusters            &  0.000027 & 0.000009 & 0.000021 & 0.000016 \\
CSCTightHaloFilter                 & 10.263068 & 0.398497 & 0.402936 & 0.508025 \\
\hline
Total                              & 28.501208 & 0.598417 & 0.601999 & 0.689380 \\
\hline
\end{tabular}

\caption{Percentage of events in data failing each of the event quality filters after requesting on good vertex. For the parked analysis we use prompt A and parked B, C and D datasets.}
\label{table_EventQualityFilterEff}
\end{table}

It is observed that the HBHENoiseFilter vetos around 22\% of the event for prompt run A but this behaviour is not observed in any of the parked datasets. This specific filter removed events with noise in the HCAL. An example of noise would be a ``hot tower'' reading very high energy values, some times for several events. Since this energy will not be balance in each event this will greatly increase \MET. On the prompt dataset there are \MET only triggers which were fired in such events which then get removed by this filter, those triggers are not present in parked dataset. This can clearly be seen if we apply first our trigger selection (Jets+\MET) and then recalculate the event rejection efficiencies, this values can be found in table \ref{table_EventQualityFilterEff_PostTrigger}.

\begin{table}[htp]
\centering

\begin{tabular}{|l||c|c|c|c|}
\hline
Filter                             & Prompt A & Parked B & Parked C & Parked D \\
\hline \hline
HBHENoiseFilter                    & 1.264571 & 0.186895 & 0.178422 & 0.154080 \\
EcalDeadCellTriggerPrimitiveFilter & 0.572468 & 0.010081 & 0.010338 & 0.010725 \\
eeBadScFilter                      & 0.000989 & 0.000001 & 0.000000 & 0.000000 \\
trackingFailureFilter              & 0.062289 & 0.000401 & 0.009729 & 0.000231 \\
manystripclus53X                   & 0.002966 & 0.001264 & 0.002304 & 0.001255 \\
toomanystripclus53X                & 0.002966 & 0.001093 & 0.001955 & 0.001103 \\
logErrorTooManyClusters            & 0.000000 & 0.000000 & 0.000000 & 0.000000 \\
CSCTightHaloFilter                 & 0.400431 & 0.397531 & 0.403156 & 0.506917 \\
\hline                             
Total                              & 2.201877 & 0.594332 & 0.592977 & 0.671026 \\
\hline
\end{tabular}

\caption{Percentage of events in data failing each of the event quality filters after requesting on good vertex and trigger conditions. For the parked analysis we use prompt A and parked B, C and D datasets.}
\label{table_EventQualityFilterEff_PostTrigger}
\end{table}

Both tables \ref{table_EventQualityFilterEff} and \ref{table_CaloLaserFilterEff} can be directly compared with the values produced by the main analysis and presented in AN-14-243\cite{CMS_AN_2014-243}. No differences are observed.

%STATUS: In progress
%TODO:  
\subsection{Jets}

In this analysis we use particle flow  jets clustered with the $anti-k_{T}$ algorithm with a cone size of 0.5. 

The correction L1FastJet, L2Relative and L3Absolute are applied to both data and monte carlo (MC) and additionally we apply L2L3Residual to data.

The to correct Jet Energy Scale we apply global tag FT53\_V21A\_AN6::All for data and START53\_V27::All for monte carlo.

The jets with this characteristics and corrections are available in the ntuples produced by the main analysis as the "standard" jets collection. We further require that our selected jets pass PFJet ID and pileup ID. And we clean the jet collection by removing all jets that closer then $\Delta R < 0.5$ to any veto electron or loose muon (relevant for control regions).

%Out of this jets we select the ones passing loose PFJet ID and pileup ID. We also clean the jet collection by removing all jets that closer then %$\Delta R < 0.5$% to any veto electron or loose muon (relevant for control regions).
 
%% FOR MY PHD THESIS 
%% Details on 
%  
% As a selection condition for out jets we use the PFJet ID described at [[CMS/JetID.Recommendations_for_8_TeV_data_a][JetID page]].
% 
% The jet requirements are:
%    * Neutral Hadron Fraction %$ <0.99 $%
%    * Neutral EM Fraction %$ <0.99 $%
%    * Number of Constituents  %$ >1 $%
% Additionally for %$|\eta| < 2.4$% we require:
%    * Charged Hadron Fraction %$  >0 $%
%    * Charged Multiplicity %$ >0 $%
%    * Charged EM Fraction %$ 0.99 $%
% 
% ---+++ Pileup ID 
% 
% Using the prescription at the [[CMS.PileupJetID][PileupJetID page]] we reject pileup jets cutting at loose working point of a BDT with Dec2012 weights and package recojets/jetproducers with version METPU_5_3_X_v4. The computation of the MVA value per jet is made when the ntuples are produced and this values are stored within each jet. We cut on the BDT score categorizing in %$p_{\perp}$% and %$|\eta|$% as follows:
% 
% %BEGINLATEX%
% \begin{tabular}{|c|c|c|}
% \hline
% Jet $p_{\perp}$ & Jet $|\eta|$ & $BDT_{score}$ \\
% \hline \hline
% $20 < p_{\perp} \leq 30$ & $|\eta| < 2.5$ & $BDT_{score} > -0.80$ \\
% $20 < p_{\perp} \leq 30$ & $2.50 \leq |\eta| < 2.75$ & $BDT_{score} > -0.85$ \\
% $20 < p_{\perp} \leq 30$ & $2.75 \leq |\eta| < 3.00$ & $BDT_{score} > -0.84$ \\
% $20 < p_{\perp} \leq 30$ & $3.00 \leq |\eta| < 5.00$ & $BDT_{score} > -0.85$ \\
% $30 < p_{\perp}$ & $|\eta| < 2.5$ & $BDT_{score} > -0.80$ \\
% $30 < p_{\perp}$ & $2.50 \leq |\eta| < 2.75$ & $BDT_{score} > -0.74$ \\
% $30 < p_{\perp}$ & $2.75 \leq |\eta| < 3.00$ & $BDT_{score} > -0.68$ \\
% $30 < p_{\perp}$ & $3.00 \leq |\eta| < 5.00$ & $BDT_{score} > -0.77$ \\
% \hline
% \end{tabular}
% %ENDLATEX%

%STATUS: In progress
%TODO:  
\subsection{Missing transverse energy \texorpdfstring{(\MET)}{}}


%STATUS: In progress
%TODO: 
% Convert base cut based ID to a table  
\subsection{Electrons}

% Make citation for CMS.EgammaCutBasedIdentification
For this analysis we use two categories of electrons "veto electrons" and "tight electrons". Both this categories of particles are based on standard EGamma POG cut based object definition which can be found at (TODO:CITATION). Additionally, we require some additional cuts to each category (TODO:WHY).

\subsubsection{Veto electrons}
 
The "veto electrons" are defined by the base requirements of the cut based electron ID veto working point of the EGamma POG with some additional cuts on top. 
 
Requirements of the cut based electron ID veto working point:

Barrel Cuts ( $ |\eta_{supercluster}|<=1.479 $ )
\begin{itemize}
  \item $ | \Delta\eta(track,supercluster) | < 0.007 $
  \item $ | \Delta\phi(track,supercluster) | < 0.8 $
  \item $ \sigma(i\eta,i\eta) < 0.01 $
  \item $ H/E < 0.15 $
  \item $ |d_{0}(vertex)| < 0.04 $
  \item $ |d_{Z}(vertex)| < 0.2 $
  \item $ \frac{PF_{isolation}}{p_{\perp}} < 0.15 $ for $ \Delta R_{cone}=0.3 $
\end{itemize}

Endcap Cuts ( $ 1.479 < |\eta_{supercluster}| < 2.5 $ )
\begin{itemize}
  \item $ | \Delta\eta(track,supercluster) | < 0.1 $
  \item $ | \Delta\phi(track,supercluster) | < 0.7 $
  \item $ \sigma(i\eta,i\eta) < 0.03 $
  \item $ |d_{0}(vertex)| < 0.04 $
  \item $ |d_{Z}(vertex)| < 0.2 $
  \item $ \frac{PF_{isolation}}{p_{\perp}} < 0.15 $ for $ \Delta R_{cone}=0.3 $
\end{itemize}

Additional requirements for this analysis
\begin{itemize}
  \item $ p_{\perp} > 10 $ GeV
  \item $ |\eta| < 2.4 $
  \item $ Effective-Area-Corrected-Isolation < 0.15 $ (is stated in the note as additional requirement but does not look like it is)
  \item $d_{xy}<0.04$ cm (is stated in the note as additional requirement but does not look like it is)
  \item $d_{z} < 0.2 $ cm (is stated in the note as additional requirement but does not look like it is)
\end{itemize}

\subsubsection{Tight electrons}

The "tight electrons" are defined by using the base requirements of the cut based electron ID tight working point (similar to 2011 very tight WP70) of the EGamma POG with some additional cuts on top.
 
Requirements of the cut based electron ID tight working point (similar to 2011 very tight WP70):

Barrel Cuts ( $ |\eta_{supercluster}|<=1.479 $ )
\begin{itemize}
  \item $ | \Delta\eta(track,supercluster) | < 0.004 $
  \item $ | \Delta\phi(track,supercluster) | < 0.3  $
  \item $ \sigma(i\eta,i\eta) < 0.01 $
  \item $ H/E < 0.12 $
  \item $ |d_{0}(vertex)| < 0.02 $
  \item $ |d_{Z}(vertex)| < 0.1 $
  \item $ |\frac{1}{E}-\frac{1}{p}| < 0.05 $
  \item $ \frac{PF_{isolation}}{p_{\perp}} < 0.10 $ for $ \Delta R_{cone}=0.3 $
  \item Conversion rejection: vertex fit probability: 1e-6
  \item Conversion rejection: missing hits $<= 0$
\end{itemize}

Endcap Cuts ( $ 1.479 < |\eta_{supercluster}| < 2.5 $ )
\begin{itemize}
  \item $ | \Delta\eta(track,supercluster) | < 0.005 $
  \item $ | \Delta\phi(track,supercluster) | < 0.2  $
  \item $ \sigma(i\eta,i\eta) < 0.03 $
  \item $ H/E < 0.10 $
  \item $ |d_{0}(vertex)| < 0.02 $
  \item $ |d_{Z}(vertex)| < 0.1 $
  \item $ |\frac{1}{E}-\frac{1}{p}| < 0.05 $
  \item $ \frac{PF_{isolation}}{p_{\perp}} < 0.10(0.07) $ for $ p_{\perp} > 20 (p_{\perp} <= 20) $ and $ \Delta R_{cone}=0.3 $
  \item Conversion rejection: vertex fit probability: 1e-6
  \item Conversion rejection: missing hits $<= 0$
\end{itemize}

Additional requirements for this analysis:
\begin{itemize}
  \item $ p_{\perp} > 20 $ GeV
  \item $ |\eta| < 2.4 $
  \item $ Effective-Area-Corrected-Isolation < 0.10 $
  \item $d_{xy}<0.02 $ cm
  \item $d_{z} < 0.1 $ cm
\end{itemize}


%STATUS: In progress
%TODO: 
\subsection{Muons}




%STATUS: In progress
%TODO: 
\subsection{Taus}


% \section{Signal selection}

\subsection{Pre-selection}

\begin{table}[!htp]
\centering

\begin{tabular}{|l||c|c|c|c|}
\hline
Filter                  & Prompt A & Parked B & Parked C & Parked D \\
\hline \hline
ECAL Laser Filter       & 0.928521 & 0.008659 & 0.000000 & 0.000000 \\
HCAL Laser Filter       & 0.007258 & 0.000000 & 0.000270 & 0.000000 \\
\hline
ECAL+ HCAL Laser Filter & 0.935704 & 0.008659 & 0.000270 & 0.000000 \\
\hline
\end{tabular}

\caption{Event rejection efficiency for the ECAL and HCAL Laser Filters. This events have to be removed due to the untimely firing of the calibration laser for this systems.}
\label{table_CaloLaserFilterEff}
\end{table}


\begin{table}[htp]
\centering

\begin{tabular}{|l|c|c|c|c||c|}
\hline
 & \rotatebox{90}{Prompt Run A} & \rotatebox{90}{Parked Run B} & \rotatebox{90}{Parked Run C} & \rotatebox{90}{Parked Run D} & \rotatebox{90}{Total Data} \\
\hline \hline
Vertex Filter & 3606391 & 132346320 & 228049748 & 308041846 & 672044305 \\
Event Quality Filters & 2658960 & 131554431 & 226680352 & 305918529 & 666812272 \\
ECAL Laser Filter & 2634271 & 131543040 & 226680352 & 305918529 & 666776192 \\
HCAL Laser Filter & 2634080 & 131543040 & 226679741 & 305918529 & 666775390 \\
L1T ETM Filter & 2461217 & 88174347 & 160560859 & 227801622 & 478998045 \\
HLT Filter & 97522 & 75100422 & 137527238 & 152041761 & 364766943 \\
$N(Electrons_{veto})=0$ & 96600 & 74947192 & 137241812 & 151725585 & 364011189 \\
$N(Muon_{loose})=0$ & 94864 & 74913002 & 137179173 & 151652654 & 363839693 \\
Dijet cut & 28164 & 23666926 & 43292391 & 42218637 & 109206118 \\
MET cut & 6252 & 57929 & 102384 & 120600 & 287165 \\
$MET_{Significance}$ cut & 3828 & 24179 & 42683 & 41620 & 112310 \\
$Min(\Delta\phi(MET,jets))$ cut & 405 & 1824 & 3452 & 3374 & 9055 \\
\hline
\end{tabular}
\caption{Event Yield for the Pre-Selection Region.}
\end{table}


\begin{table}[!htp]
\centering

\begin{tabular}{|l|c|c|c|}
\hline
Dataset & Main Analysis & Cross Check Analysis & $\frac{CC}{Main}-1$ \\ 
\hline \hline
Prompt A &  405 &  405 & 0.00\% \\
Parked B & 1824 & 1824 & 0.00\% \\
Parked C & 3453 & 3452 & -0.03 \% \\
Parked D & 3374 & 3374 & 0.0\% \\
\hline \hline
Total & 9056 & 9055 & -0.01\% \\
\hline
\end{tabular}

\caption{Comparison between main and cross check analysis for the event yield of the event pre-selection. There is a difference of a single event between both analysis which is a difference in total yield around 0.01\%.}
\end{table}


\subsection{Signal region}

\begin{table}[htp]
\centering

\begin{tabular}{|l|c|c|c|c||c|}
\hline
 & \rotatebox{90}{Prompt Run A} & \rotatebox{90}{Parked Run B} & \rotatebox{90}{Parked Run C} & \rotatebox{90}{Parked Run D} & \rotatebox{90}{Total Data} \\
\hline \hline
Vertex Filter & 3606391 & 132346320 & 228049748 & 308041846 & 672044305 \\
Event Quality Filters & 2658960 & 131554431 & 226680352 & 305918529 & 666812272 \\
ECAL Laser Filter & 2634271 & 131543040 & 226680352 & 305918529 & 666776192 \\
HCAL Laser Filter & 2634080 & 131543040 & 226679741 & 305918529 & 666775390 \\
L1T ETM Filter & 2461217 & 88174347 & 160560859 & 227801622 & 478998045 \\
HLT Filter & 97522 & 75100422 & 137527238 & 152041761 & 364766943 \\
$N(Electrons_{veto})=0$ & 96600 & 74947192 & 137241812 & 151725585 & 364011189 \\
$N(Muon_{loose})=0$ & 94864 & 74913002 & 137179173 & 151652654 & 363839693 \\
Dijet cut & 18338 & 13678405 & 25090291 & 24082304 & 62869338 \\
MET cut & 4167 & 38178 & 68047 & 79723 & 190115 \\
$MET_{Significance}$ cut & 786 & 3396 & 5988 & 5567 & 15737 \\
$Min(\Delta\phi(MET,jets))$ cut & 34 & 91 & 205 & 178 & 508 \\
\hline
\end{tabular}
\caption{Event Yield for the Signal Region.}
\end{table}


\begin{table}[!htp]
\centering

\begin{tabular}{|l|c|c|c|}
\hline
Dataset & Main Analysis & Cross Check Analysis & $\frac{CC}{Main}-1$ \\ 
\hline \hline
Prompt A &  34 &  34 & 0.00\% \\
Parked B &  91 &  91 & 0.00\% \\
Parked C & 205 & 205 & 0.00 \% \\
Parked D & 178 & 178 & 0.00\% \\
\hline \hline
Total & 508 & 508 & 0.00\% \\
\hline
\end{tabular}

\caption{Comparison between main and cross check analysis for the event yield of signal region. No difference in yields is observed either in total or by acquisition era.}
\end{table}

% \section{Background estimation}

\subsection{W to \texorpdfstring{electron+\MET}{electron+MET}}

The selection consists of the following cuts:

\begin{itemize}
  \item Vertex cut
  \item Event quality filters (MET Filters)
  \item ECAL+HCAL Laser filters
  \item L1T ETM Filter (L1T\_ETM40 emulation)
  \begin{itemize}
    \item $ L1T\_ETM >= 40 $
  \end{itemize}
  \item HLT path filter
  \begin{itemize}
    \item From run 190456 to 193621 (Run 2012 A) use HLT\_DiPFJet40\_PFMETnoMu65\_MJJ800VBF\_AllJets* 
    \item From run 193833 to 196531 (Run 2012 B) use HLT\_DiJet35\_MJJ700\_AllJets\_DEta3p5\_VBF*
    \item From run 198022 to 203742 (Run 2012 C) use HLT\_DiJet35\_MJJ700\_AllJets\_DEta3p5\_VBF*
    \item From run 203777 to 208686 (Run 2012 D) use HLT\_DiJet30\_MJJ700\_AllJets\_DEta3p5\_VBF*
  \end{itemize}
  \item Exactly one $Electron_{Veto}$
  \begin{itemize}
    \item Using veto electron defined on this page
  \end{itemize}
  \item Exactly one $Electron_{Tight}$
  \begin{itemize}
    \item Using tight electron defined on this page
  \end{itemize}
  \item Muon Veto
  \begin{itemize}
    \item Using veto muons defined on this page (to be done)
  \end{itemize}
  \item $ MET > 90 $ GeV
  \item $ MET_{significance} > 4.0 $
  \item Dijet cut (leading dijet requirements):
  \begin{itemize}
    \item Lead dijet $ p_{\perp} > 50$ GeV
    \item Sub-lead dijet $ p_{\perp} > 45$ GeV
    \item Jets $ |\eta| < 4.7 $
    \item Dijet $ \Delta\eta < 3.6 $
    \item Dijet $ m_{jj} > 1200 $ GeV
  \end{itemize}
  \item $ Min(\Delta\phi(MET,Jet_{p_{\perp}>30})))>2.3 $
\end{itemize}

\begin{table}[!htp]
\centering

\begin{tabular}{|l|c|c|c|c||c|}
\hline
 & \rotatebox{90}{Prompt Run A} & \rotatebox{90}{Parked Run B} & \rotatebox{90}{Parked Run C} & \rotatebox{90}{Parked Run D} & \rotatebox{90}{Total Data} \\
\hline \hline
Vertex Filter & 3606391 & 132346320 & 228049748 & 308041846 & 672044305 \\
Event Quality Filters & 2658960 & 131554431 & 226680352 & 305918529 & 666812272 \\
ECAL Laser Filter & 2634271 & 131543040 & 226680352 & 305918529 & 666776192 \\
HCAL Laser Filter & 2634080 & 131543040 & 226679741 & 305918529 & 666775390 \\
L1T ETM Filter & 2461217 & 88174347 & 160560859 & 227801622 & 478998045 \\
HLT Filter & 97522 & 75100422 & 137527238 & 152041761 & 364766943 \\
$N(Electrons_{veto})=1$ & 899 & 151621 & 282481 & 312807 & 747808 \\
$N(Muon_{loose})=0$ & 852 & 151249 & 281801 & 312033 & 745935 \\
$N(Electrons_{tight})=1$ & 398 & 23751 & 44323 & 50279 & 118751 \\
Dijet cut & 64 & 1607 & 3145 & 3077 & 7893 \\
MET cut & 55 & 281 & 543 & 525 & 1404 \\
$MET_{Significance}$ cut & 31 & 123 & 242 & 197 & 593 \\
$Min(\Delta\phi(MET,jets))$ cut & 4 & 16 & 24 & 24 & 68 \\
\hline
\end{tabular}

\caption{Electron+MET Yields}
\end{table}


\begin{table}[!htp]
  \centering
  
\begin{tabular}{|l|c|c||c|}
  \hline
  Dataset & Main Analysis & Cross Check Analysis & $\frac{CC}{Main}-1$ \\ 
  \hline \hline
  Prompt Run A &  4 &  4 & 0.00\% \\
  Parked Run B & 16 & 16 & 0.00\% \\
  Parked Run C & 24 & 24 & 0.00\% \\
  Parked Run D & 24 & 24 & 0.00\% \\
  \hline \hline
  Total & 68 & 68 & 0.00\% \\
  \hline
\end{tabular}

\caption{Comparison between main and cross check analysis for the event yield of Electron+MET event selection. No difference in yields is observed either in total or by acquisition era.}
\end{table}

\subsection{W to \texorpdfstring{$\mu$+\MET}{muon+MET}}

The selection consists of the following cuts
\begin{itemize}
  \item Vertex cut 
  \item Event quality filters (MET Filters)
  \item ECAL+HCAL Laser filters
  \item L1T ETM Filter (L1T\_ETM40 emulation)
  \begin{itemize}
    \item $ L1T\_ETM >= 40 $
  \end{itemize}
  \item HLT path filter
  \begin{itemize}
    \item From run 190456 to 193621 (Run 2012 A) use HLT\_DiPFJet40\_PFMETnoMu65\_MJJ800VBF\_AllJets* 
    \item From run 193833 to 196531 (Run 2012 B) use HLT\_DiJet35\_MJJ700\_AllJets\_DEta3p5\_VBF*
    \item From run 198022 to 203742 (Run 2012 C) use HLT\_DiJet35\_MJJ700\_AllJets\_DEta3p5\_VBF*
    \item From run 203777 to 208686 (Run 2012 D) use HLT\_DiJet30\_MJJ700\_AllJets\_DEta3p5\_VBF*
  \end{itemize}
  \item Electron Veto
  \begin{itemize}
    \item Using veto electron defined on this page
  \end{itemize}
  \item Exactly one $Muon_{loose}$
  \begin{itemize}
    \item Using veto muons defined on this page (to be done)
  \end{itemize}
  \item Exactly one $Muon_{tight}$
  \begin{itemize}
    \item Using veto muons defined on this page (to be done)
  \end{itemize}
  \item $MET > 90 $ GeV
  \item $MET_{significance} > 4.0 $
  \item Dijet cut (leading dijet requirements):
  \begin{itemize}
    \item Lead dijet $ p_{\perp} > 50$ GeV
    \item Sub-lead dijet $ p_{\perp} > 45$ GeV
    \item Jets $ |\eta| < 4.7 $
    \item Dijet $ \Delta\eta < 3.6 $
    \item Dijet $ m_{jj} > 1200 $ GeV
  \end{itemize}
  \item $ Min(\Delta\phi(MET,Jet_{p_{\perp}>30})))>2.3 $
\end{itemize}

\begin{table}[!htp]
\centering

\begin{tabular}{|l|c|c|c|c||c|}
\hline
 & \rotatebox{90}{Prompt Run A} & \rotatebox{90}{Parked Run B} & \rotatebox{90}{Parked Run C} & \rotatebox{90}{Parked Run D} & \rotatebox{90}{Total Data} \\
\hline \hline
Vertex Filter & 3606391 & 132346320 & 228049748 & 308041846 & 672044305 \\
Event Quality Filters & 2658960 & 131554431 & 226680352 & 305918529 & 666812272 \\
ECAL Laser Filter & 2634271 & 131543040 & 226680352 & 305918529 & 666776192 \\
HCAL Laser Filter & 2634080 & 131543040 & 226679741 & 305918529 & 666775390 \\
L1T ETM Filter & 2461217 & 88174347 & 160560859 & 227801622 & 478998045 \\
HLT Filter & 97522 & 75100422 & 137527238 & 152041761 & 364766943 \\
$N(Electrons_{veto})=0$ & 96600 & 74947192 & 137241812 & 151725585 & 364011189 \\
$N(Muon_{loose})=1$ & 1625 & 33505 & 61325 & 71449 & 167904 \\
$N(Muon_{tight})=1$ & 1223 & 9662 & 17873 & 20088 & 48846 \\
Dijet cut & 176 & 1493 & 2740 & 2684 & 7093 \\
MET cut & 157 & 809 & 1545 & 1493 & 4004 \\
$MET_{Significance}$ cut & 86 & 487 & 910 & 825 & 2308 \\
$Min(\Delta\phi(MET,jets))$ cut & 10 & 60 & 124 & 106 & 300 \\
\hline
\end{tabular}
\caption{Muon+MET Region}
\end{table}


\begin{table}[!htp]
\centering

\begin{tabular}{|l|c|c||c|}
  \hline
  Dataset & Main Analysis & Cross Check Analysis & $\frac{CC}{Main}-1$ \\ 
  \hline \hline
  Prompt Run A &  10 &  10 & 0.00\% \\
  Parked Run B &  60 &  60 & 0.00\% \\
  Parked Run C & 124 & 124 & 0.00\% \\
  Parked Run D & 106 & 106 & 0.00\% \\
  \hline \hline
  Total & 300 & 300 & 0.00\% \\
  \hline
\end{tabular}

\caption{Comparison between main and cross check analysis for the event yield of Muon+MET event selection. No difference in yields is observed either in total or by acquisition era.}
\end{table}

\subsection{W to \texorpdfstring{$\tau$+\MET}{tau+MET}}

\begin{table}[!htp]
\centering

\begin{tabular}{|l|c|c|c|c||c|}
\hline
 & \rotatebox{90}{Prompt Run A} & \rotatebox{90}{Parked Run B} & \rotatebox{90}{Parked Run C} & \rotatebox{90}{Parked Run D} & \rotatebox{90}{Total Data} \\
\hline \hline
Vertex Filter & 3606391 & 132346320 & 228049748 & 308041846 & 672044305 \\
Event Quality Filters & 2658960 & 131554431 & 226680352 & 305918529 & 666812272 \\
ECAL Laser Filter & 2634271 & 131543040 & 226680352 & 305918529 & 666776192 \\
HCAL Laser Filter & 2634080 & 131543040 & 226679741 & 305918529 & 666775390 \\
L1T ETM Filter & 2461217 & 88174347 & 160560859 & 227801622 & 478998045 \\
HLT Filter & 97522 & 75100422 & 137527238 & 152041761 & 364766943 \\
$N(Electrons_{veto})=0$ & 96600 & 74947192 & 137241812 & 151725585 & 364011189 \\
$N(Muon_{loose})=0$ & 94864 & 74913002 & 137179173 & 151652654 & 363839693 \\
Dijet cut & 18338 & 13678405 & 25090291 & 24082304 & 62869338 \\
MET cut & 4167 & 38178 & 68047 & 79723 & 190115 \\
$MET_{Significance}$ cut & 786 & 3396 & 5988 & 5567 & 15737 \\
$N(Tau)=1$ & 12 & 47 & 63 & 59 & 181 \\
$M_{perp}(MET,\tau)$ cut & 5 & 35 & 46 & 38 & 124 \\
$Min(\Delta\phi(MET,Dijet))$ cut & 2 & 22 & 25 & 27 & 76 \\
\hline
\end{tabular}
\caption{Event Yield for the Tau+MET Region.}
\end{table}


\begin{table}[!htp]
\centering

\begin{tabular}{|l|c|c||c|}
  \hline
  Dataset & Main Analysis & Cross Check Analysis & $\frac{CC}{Main}-1$ \\ 
  \hline \hline
  Prompt Run A &  2 &  2 & 0.00\% \\
  Parked Run B & 22 & 22 & 0.00\% \\
  Parked Run C & 25 & 25 & 0.00\% \\
  Parked Run D & 27 & 27 & 0.00\% \\
  \hline \hline
  Total & 76 & 76 & 0.00\% \\
  \hline
\end{tabular}

\caption{Comparison between main and cross check analysis for the event yield of Tau+MET event selection. No difference in yields is observed either in total or by acquisition era.}
\end{table}

\subsection{Z to \texorpdfstring{$\mu\mu$}{mumu}}

The selection consists of the following cuts:
\begin{itemize}
  \item Vertex cut
  \item Event quality filters (MET Filters)
  \item ECAL+HCAL Laser filters
  \item !L1T ETM Filter (L1T\_ETM40 emulation)
  \begin{itemize}
    \item $ L1T\_ETM >= 40 $
  \end{itemize}
  \item HLT path filter
  \begin{itemize}
    \item From run 190456 to 193621 (Run 2012 A) use HLT\_DiPFJet40\_PFMETnoMu65\_MJJ800VBF\_AllJets* 
    \item From run 193833 to 196531 (Run 2012 B) use HLT\_DiJet35\_MJJ700\_AllJets\_DEta3p5\_VBF*
    \item From run 198022 to 203742 (Run 2012 C) use HLT\_DiJet35\_MJJ700\_AllJets\_DEta3p5\_VBF*
    \item From run 203777 to 208686 (Run 2012 D) use HLT\_DiJet30\_MJJ700\_AllJets\_DEta3p5\_VBF*
  \end{itemize}
  \item Electron Veto
  \begin{itemize}
    \item Using veto electron defined on this page
  \end{itemize}
  \item Exactly two $Muon_{loose}$
  \begin{itemize}
    \item Using veto muons defined on this page (to be done)
  \end{itemize}
  \item Exactly two $Muon_{tight}$
  \begin{itemize}
    \item Using veto muons defined on this page (to be done)
  \end{itemize}
  \item Dimuon with $60<mass<120$ GeV
  \item $ MET > 90 $ GeV
  \item $ MET_{significance} > 4.0 $
  \item Dijet cut (leading dijet requirements):
  \begin{itemize}
    \item Lead dijet $ p_{\perp} > 50$ GeV
    \item Sub-lead dijet $ p_{\perp} > 45$ GeV
    \item Jets $ |\eta| < 4.7 $
    \item Dijet $ \Delta\eta < 3.6 $
    \item Dijet $ m_{jj} > 1200 $ GeV
  \end{itemize}
\end{itemize}
  
\begin{table}[!htp]
\centering

\begin{tabular}{|l|c|c|c|c||c|}
\hline
 & \rotatebox{90}{Prompt Run A} & \rotatebox{90}{Parked Run B} & \rotatebox{90}{Parked Run C} & \rotatebox{90}{Parked Run D} & \rotatebox{90}{Total Data} \\
\hline \hline
Vertex Filter & 3606391 & 132346320 & 228049748 & 308041846 & 672044305 \\
Event Quality Filters & 2658960 & 131554431 & 226680352 & 305918529 & 666812272 \\
ECAL Laser Filter & 2634271 & 131543040 & 226680352 & 305918529 & 666776192 \\
HCAL Laser Filter & 2634080 & 131543040 & 226679741 & 305918529 & 666775390 \\
L1T ETM Filter & 2461217 & 88174347 & 160560859 & 227801622 & 478998045 \\
HLT Filter & 97522 & 75100422 & 137527238 & 152041761 & 364766943 \\
$N(Electrons_{veto})=0$ & 96600 & 74947192 & 137241812 & 151725585 & 364011189 \\
selTwoMuonsLoose & 111 & 683 & 1312 & 1480 & 3586 \\
selTwoMuonsTight & 73 & 450 & 822 & 935 & 2280 \\
Z Mass cut & 62 & 379 & 686 & 768 & 1895 \\
Dijet cut & 11 & 58 & 98 & 96 & 263 \\
MET cut & 9 & 41 & 74 & 70 & 194 \\
$MET_{Significance}$ cut & 7 & 18 & 55 & 44 & 124 \\
$Min(\Delta\phi(MET,jets))$ cut & 2 & 4 & 5 & 7 & 18 \\
\hline
\end{tabular}
\caption{Event Yield for the Z to $\mu\mu$ Region.}
\end{table}


\begin{table}[!htp]
\centering

\begin{tabular}{|l|c|c||c|}
  \hline
  Dataset & Main Analysis & Cross Check Analysis & $\frac{CC}{Main}-1$ \\ 
  \hline \hline
  Prompt A & 2 & 2 & 0.00\% \\
  Parked B & 4 & 4 & 0.00\% \\
  Parked C & 5 & 5 & 0.00\% \\
  Parked D & 7 & 7 & 0.00\% \\
  \hline \hline
  Total & 18 & 18 & 0.00\% \\
  \hline
\end{tabular}

\caption{Comparison between main and cross check analysis for the event yield of Z to $\mu\mu$ event selection. No difference in yields is observed either in total or by acquisition era.}
\end{table}

\subsection{Top}

The selection consists of the following cuts:
\begin{itemize}
  \item Vertex cut 
  \item Event quality filters (MET Filters)
  \item ECAL+HCAL Laser filters
  \item L1T ETM Filter (L1T\_ETM40 emulation)
  \begin{itemize}
    \item $ L1T\_ETM >= 40 $
  \end{itemize}
  \item HLT path filter
  \begin{itemize}
    \item From run 190456 to 193621 (Run 2012 A) use HLT\_DiPFJet40\_PFMETnoMu65\_MJJ800VBF\_AllJets* 
    \item From run 193833 to 196531 (Run 2012 B) use HLT\_DiJet35\_MJJ700\_AllJets\_DEta3p5\_VBF*
    \item From run 198022 to 203742 (Run 2012 C) use HLT\_DiJet35\_MJJ700\_AllJets\_DEta3p5\_VBF*
    \item From run 203777 to 208686 (Run 2012 D) use HLT\_DiJet30\_MJJ700\_AllJets\_DEta3p5\_VBF*
  \end{itemize}
  \item Exactly one $Electron_{Veto}$
  \begin{itemize}
    \item Using veto electron defined on this page
  \end{itemize}
  \item Exactly one $Electron_{Tight}$
  \begin{itemize}
    \item Using tight electron defined on this page
  \end{itemize}
  \item Exactly one $Muon_{loose}$
  \begin{itemize}
    \item Using veto muons defined on this page (to be done)
  \end{itemize}
  \item Exactly one $Muon_{tight}$
  \begin{itemize}
    \item Using veto muons defined on this page (to be done)
  \end{itemize}
  \item $MET > 90 $ GeV
  \item $MET_{significance} > 4.0 $
  \item Dijet cut (leading dijet requirements):
  \begin{itemize}
    \item Lead dijet $ p_{\perp} > 50$ GeV
    \item Sub-lead dijet $ p_{\perp} > 45$ GeV
    \item Jets $ |\eta| < 4.7 $
    \item Dijet $ \Delta\eta < 3.6 $
    \item Dijet $ m_{jj} > 1200 $ GeV
  \end{itemize}
\end{itemize}
    
\begin{table}[!htp]
\centering

\begin{tabular}{|l|c|c|c|c||c|}
\hline
 & \rotatebox{90}{Prompt Run A} & \rotatebox{90}{Parked Run B} & \rotatebox{90}{Parked Run C} & \rotatebox{90}{Parked Run D} & \rotatebox{90}{Total Data} \\
\hline \hline
Vertex Filter & 3606391 & 132346320 & 228049748 & 308041846 & 672044305 \\
Event Quality Filters & 2658960 & 131554431 & 226680352 & 305918529 & 666812272 \\
ECAL Laser Filter & 2634271 & 131543040 & 226680352 & 305918529 & 666776192 \\
HCAL Laser Filter & 2634080 & 131543040 & 226679741 & 305918529 & 666775390 \\
L1T ETM Filter & 2461217 & 88174347 & 160560859 & 227801622 & 478998045 \\
HLT Filter & 97522 & 75100422 & 137527238 & 152041761 & 364766943 \\
$N(Muon_{loose})=1$ & 1674 & 33877 & 61994 & 72215 & 169760 \\
$N(Muon_{tight})=1$ & 1259 & 9913 & 18322 & 20585 & 50079 \\
$N(Electrons_{veto})=1$ & 35 & 247 & 444 & 492 & 1218 \\
$N(Electrons_{tight})=1$ & 23 & 150 & 273 & 300 & 746 \\
Dijet cut & 0 & 15 & 21 & 17 & 53 \\
MET cut & 0 & 10 & 14 & 12 & 36 \\
$MET_{Significance}$ cut & 0 & 4 & 9 & 8 & 21 \\
\hline
\end{tabular}
\caption{Event Yield for the Top Region.}
\end{table}


\begin{table}[!htp]
\centering

\begin{tabular}{|l|c|c||c|}
  \hline
  Dataset & Main Analysis & Cross Check Analysis & $\frac{CC}{Main}-1$ \\ 
  \hline \hline
  Prompt A & 0 & 0 & 0.00\% \\
  Parked B & 4 & 4 & 0.00\% \\
  Parked C & 9 & 9 & 0.00\% \\
  Parked D & 8 & 8 & 0.00\% \\
  \hline \hline
  Total & 21 & 21 & 0.00\% \\
  \hline
\end{tabular}

\caption{Comparison between main and cross check analysis for the event yield of top event selection. No difference in yields is observed either in total or by acquisition era.}
\end{table}

\subsection{QCD}

\begin{table}[htp]
\centering

\begin{tabular}{|l|c|c|c|c||c|}
\hline
 & \rotatebox{90}{Prompt Run A} & \rotatebox{90}{Parked Run B} & \rotatebox{90}{Parked Run C} & \rotatebox{90}{Parked Run D} & \rotatebox{90}{Total Data} \\
\hline \hline
Vertex Filter & 3606391 & 132346320 & 228049748 & 308041846 & 672044305 \\
Event Quality Filters & 2658960 & 131554431 & 226680352 & 305918529 & 666812272 \\
ECAL Laser Filter & 2634271 & 131543040 & 226680352 & 305918529 & 666776192 \\
HCAL Laser Filter & 2634080 & 131543040 & 226679741 & 305918529 & 666775390 \\
L1T ETM Filter & 2461217 & 88174347 & 160560859 & 227801622 & 478998045 \\
HLT Filter & 97522 & 75100422 & 137527238 & 152041761 & 364766943 \\
$N(Electrons_{veto})=0$ & 96600 & 74947192 & 137241812 & 151725585 & 364011189 \\
$N(Muon_{loose})=0$ & 94864 & 74913002 & 137179173 & 151652654 & 363839693 \\
Dijet cut & 18338 & 13678405 & 25090291 & 24082304 & 62869338 \\
MET cut & 4167 & 38178 & 68047 & 79723 & 190115 \\
$MET_{Significance}$ cut & 2532 & 15594 & 27623 & 27068 & 72817 \\
$Min(\Delta\phi(MET,jets))$ cut & 2314 & 14691 & 25826 & 25326 & 68157 \\
\hline
\end{tabular}
\caption{Event Yield for the QCD Region.}
\end{table}





  \chapter{Technical work}

\section{Trigger work}
  \chapter{Physics Analysis}

\section{Prompt/Parked trigger studies}

\section{Prompt Analysis}

\section{Parked Analysis}

\section{Run II trigger studies}

\section{Run II Analysis}


  \chapter{Conclusions}

Summary of relevant results and their impact on Particle Physics
\end{mainmatter}

\begin{backmatter}
  \begin{colophon}
  This thesis was made in \LaTeXe{} using the ``hepthesis'' class~\cite{hepthesis}.
\end{colophon}

%% You're recommended to use the eprint-aware biblio styles which
%% can be obtained from e.g. www.arxiv.org. The file mythesis.bib
%% is derived from the source using the SPIRES Bibtex service.
\bibliographystyle{h-physrev}
\bibliography{mythesis}

%% I prefer to put these tables here rather than making the
%% front matter seemingly interminable. No-one cares, anyway!
\listoffigures
\listoftables

%% If you have time and interest to generate a (decent) index,
%% then you've clearly spent more time on the write-up than the 
%% research ;-)
%\printindex

\end{backmatter}

\end{document}
